\documentclass[11pt]{article}
\usepackage{amsmath,amssymb,amsthm}
\usepackage{algorithm}
\usepackage{algpseudocode}
\usepackage{hyperref}
\usepackage{enumitem}
\usepackage{bm}
\usepackage{geometry}
\geometry{margin=1in}
\newtheorem{theorem}{Theorem}
\newtheorem{lemma}{Lemma}
\newtheorem{assumption}{Assumption}
\newtheorem{corollary}{Corollary}
\newcommand{\E}{\mathbb{E}}
\newcommand{\R}{\mathbb{R}}

\begin{document}

\section*{RMSProp for Non‑Convex Smooth Optimization}

\section*{Mathematical Formulation}
Let $f:\R^{d}\rightarrow \R$ be differentiable.  
At iteration $t\ge 1$ we have a parameter vector $w_t\in\R^{d}$ and a stochastic gradient 
\[
g_t = \nabla f(w_t;\xi_t),
\]
where $\xi_t$ denotes the random data sample (or mini‑batch) drawn at iteration $t$.

The RMSProp algorithm maintains an exponential moving average of the element‑wise squared gradients:
\begin{align}
    v_t &= \beta\, v_{t-1} + (1-\beta)\, g_t\odot g_t,
    \label{eq:vt}
\end{align}
where $v_t\in\R^{d}$, $0<\beta<1$ and $\odot$ denotes element‑wise product.  
The learning rate for each coordinate is adapted by the inverse square root of $v_t$:
\begin{align}
    w_{t+1} &= w_t - \alpha \, \frac{g_t}{\sqrt{v_t}+\epsilon},
    \label{eq:update}
\end{align}
with step‑size hyperparameter $\alpha>0$ and a small constant $\epsilon>0$ for numerical stability.

\section*{Pseudocode}
\begin{algorithm}[H]
\caption{RMSProp (non‑convex setting)}\label{alg:rmsprop}
\begin{algorithmic}[1]
\State \textbf{Input:} initial point $w_0\in\R^{d}$, step size $\alpha>0$, decay $\beta\in(0,1)$, $\epsilon>0$, total iterations $T$.
\State Initialize $v_0\gets \mathbf{0}\in\R^{d}$.
\For{$t=0,\dots,T-1$}
    \State Sample $\xi_t$ and compute stochastic gradient $g_t=\nabla f(w_t;\xi_t)$.
    \State $v_{t+1}\gets \beta\,v_t+(1-\beta)\,g_t\odot g_t$ \hfill\Comment{Eq.~\eqref{eq:vt}}
    \State $w_{t+1}\gets w_t-\alpha\,\dfrac{g_t}{\sqrt{v_{t+1}}+\epsilon}$ \hfill\Comment{Eq.~\eqref{eq:update}}
\EndFor
\State \textbf{Output:} $w_T$ (or the average $\bar w_T=\frac1T\sum_{t=0}^{T-1}w_t$).
\end{algorithmic}
\end{algorithm}

\section*{Dry Run Example}
Consider the one‑dimensional quadratic $f(w)=(w-3)^2$.  
Its exact gradient is $\nabla f(w)=2(w-3)$.  
We use a deterministic ``stochastic’’ gradient (i.e. $g_t=\nabla f(w_t)$) with parameters
\[
\alpha=0.1,\qquad \beta=0.9,\qquad \epsilon=10^{-8},\qquad w_0=0.
\]

\begin{center}
\begin{tabular}{c|c|c|c|c}
$t$ & $w_t$ & $g_t=2(w_t-3)$ & $v_t$ (recurrence) & $w_{t+1}=w_t-\alpha g_t/(\sqrt{v_t}+\epsilon)$\\\hline
0 & $0$ & $-6$ & $v_1=0.9\cdot0+(1-0.9)6^2=0.1\cdot36=3.6$ &
$0-0.1\cdot(-6)/\sqrt{3.6}=0+0.6/1.897\approx0.316$\\\hline
1 & $0.316$ & $2(0.316-3)=-5.368$ &
$v_2=0.9\cdot3.6+(0.1)5.368^2=3.24+0.1\cdot28.82\approx6.122$ &
$0.316-0.1\cdot(-5.368)/\sqrt{6.122}=0.316+0.5368/2.474\approx0.537$\\\hline
2 & $0.537$ & $2(0.537-3)=-4.926$ &
$v_3=0.9\cdot6.122+0.1\cdot4.926^2=5.5098+0.1\cdot24.27\approx7.937$ &
$0.537-0.1\cdot(-4.926)/\sqrt{7.937}=0.537+0.4926/2.819\approx0.714$\\\hline
\end{tabular}
\end{center}
The objective values $f(w_t)$ decrease from $9$ (at $w_0$) to $6.28$, $5.41$, $4.61$ after three iterations, illustrating the adaptive step‑size effect.

\newpage
\section*{Assumptions}
\begin{assumption}[Smoothness]\label{ass:smooth}
$f$ is $L$‑smooth, i.e.\ for all $x,y\in\R^{d}$,
\[
\|\nabla f(x)-\nabla f(y)\|\le L\|x-y\|.
\]
Equivalently,
\[
f(y)\le f(x)+\langle\nabla f(x),y-x\rangle+\frac{L}{2}\|y-x\|^2.
\]
\end{assumption}
\begin{assumption}[Unbiased Stochastic Gradient]\label{ass:unbiased}
For every $t$,
\[
\E\!\left[g_t\mid w_t\right]=\nabla f(w_t),
\]
and the conditional second moment is bounded:
\[
\E\!\left[\|g_t\|^2\mid w_t\right]\le G^2,
\]
for some $G>0$.
\end{assumption}
\begin{assumption}[Bounded Variance]\label{ass:var}
There exists $\sigma^2\ge0$ such that
\[
\E\!\left[\|g_t-\nabla f(w_t)\|^2\mid w_t\right]\le\sigma^2.
\]
\end{assumption}
\begin{assumption}[Non‑zero Denominator]\label{ass:epsilon}
$\epsilon>0$ is chosen so that $\sqrt{v_{t,i}}+\epsilon\ge \epsilon$ for all coordinates $i$ and all $t$.
\end{assumption}
All constants $\alpha,\beta,\epsilon$ are fixed throughout the run.

\section*{Main Convergence Theorem}
\begin{theorem}[Convergence of RMSProp for non‑convex $L$‑smooth functions]\label{thm:main}
Under Assumptions~\ref{ass:smooth}–\ref{ass:epsilon}, let $\{w_t\}_{t=0}^{T}$ be generated by Algorithm~\ref{alg:rmsprop}.  
Define the (scalar) adaptive learning rate
\[
\eta_t \;:=\; \frac{\alpha}{\sqrt{\bar v_t}+\epsilon},
\qquad 
\bar v_t \;:=\;\frac{1}{d}\sum_{i=1}^{d} v_{t,i}.
\]
Assume the decay parameter satisfies $\beta\in[0,1)$ and let 
\[
\gamma \;:=\; \min_{t\ge0}\frac{\alpha}{\sqrt{\bar v_t}+\epsilon}>0.
\]
Then for any $T\ge1$,
\[
\frac{1}{T}\sum_{t=0}^{T-1}\E\!\big[\|\nabla f(w_t)\|^2\big]
\;\le\;
\frac{2\big(f(w_0)-f^\star\big)}{\gamma T}
\;+\;
\frac{L\alpha^2 G^2}{\gamma^2 T}
\;+\;
\frac{2L\alpha\sigma}{\gamma\sqrt{T}},
\tag{1}
\]
where $f^\star:=\inf_{w}f(w)$. Consequently,
\[
\min_{0\le t<T}\E\!\big[\|\nabla f(w_t)\|^2\big]
= O\!\Bigl(\frac{1}{\sqrt{T}}\Bigr).
\]
\end{theorem}

\section*{Theoretical Properties}
\begin{itemize}[leftmargin=*]
\item \textbf{Stability.} The exponential moving average in~\eqref{eq:vt} guarantees that $v_t$ never explodes; indeed $v_{t,i}\le\max\{v_{0,i},(1-\beta)G^2\}$ almost surely.
\item \textbf{Hyper‑parameter Sensitivity.} Larger $\beta$ yields slower adaptation (more inertia), which tightens the bound via a larger $\gamma$ but may increase the bias term $L\alpha^2G^2/(\gamma^2T)$. Conversely, a smaller $\beta$ accelerates adaptation but can reduce $\gamma$.
\item \textbf{Comparison to SGD.} Setting $v_t\equiv 1$ and $\epsilon=0$ recovers vanilla SGD; Theorem~\ref{thm:main} reduces to the classical $O(1/\sqrt{T})$ gradient‑norm bound. RMSProp gains the same asymptotic rate while often exhibiting a smaller constant because $\gamma$ adapts to the magnitude of past gradients.
\item \textbf{Special Cases.}
\begin{enumerate}
    \item If $g_t$ is deterministic (full‑batch gradient), $\sigma=0$ and the bound simplifies to $O(1/T)$ for the average gradient norm.
    \item When $L\alpha\le\gamma$, the second term in~\eqref{1} is dominated by the first, yielding a clean $O(1/\sqrt{T})$ rate.
\end{enumerate}
\end{itemize}

\section*{Discussion}
The theorem shows that RMSProp, despite its per‑coordinate adaptive step size, enjoys the same worst‑case convergence guarantee as SGD for smooth non‑convex objectives. The proof hinges on a careful handling of the moving average $v_t$ to bound the effective learning rate from below (via $\gamma$) and from above (via the smoothness constant $L$). The additional variance term $\sigma$ appears only as $O(1/\sqrt{T})$, matching the optimal stochastic rate.

\newpage
\section*{Preliminaries (Definitions and Lemmas)}
\begin{lemma}[Descent Lemma]\label{lem:descent}
If $f$ is $L$‑smooth, then for any $x,y\in\R^{d}$,
\[
f(y)\le f(x)+\langle\nabla f(x),y-x\rangle+\frac{L}{2}\|y-x\|^2.
\]
\end{lemma}
\begin{lemma}[Bound on the Adaptive Learning Rate]\label{lem:eta}
Let $\eta_t$ be defined as in Theorem~\ref{thm:main}. Then for all $t$,
\[
\gamma\;\le\;\eta_t\;\le\;\frac{\alpha}{\epsilon}.
\]
\end{lemma}
\begin{proof}
The lower bound follows directly from the definition of $\gamma$ as the minimum of $\alpha/(\sqrt{\bar v_t}+\epsilon)$.  
The upper bound uses $\sqrt{\bar v_t}\ge0$ and $\epsilon>0$.
\end{proof}

\section*{Main Proof (Step‑by‑Step Derivation)}
We prove Theorem~\ref{thm:main}. All expectations are taken w.r.t. the randomness up to time $t$ unless otherwise noted.

\begin{enumerate}
\item[Step 1.] \textbf{Apply the Descent Lemma} (Lemma~\ref{lem:descent}) with $x=w_t$, $y=w_{t+1}=w_t-\eta_t g_t$:
\begin{align}
    f(w_{t+1})
    &\le f(w_t)-\eta_t\langle\nabla f(w_t),g_t\rangle
          +\frac{L}{2}\eta_t^{2}\|g_t\|^{2}.
    \label{eq:descent}
\end{align}
\item[Step 2.] \textbf{Take conditional expectation} given $w_t$ and use unbiasedness (Assumption~\ref{ass:unbiased}):
\begin{align}
    \E\!\big[f(w_{t+1})\mid w_t\big]
    &\le f(w_t)-\eta_t\|\nabla f(w_t)\|^{2}
      +\frac{L}{2}\eta_t^{2}\E\!\big[\|g_t\|^{2}\mid w_t\big].
    \label{eq:condexp}
\end{align}
\item[Step 3.] \textbf{Bound the second moment} of $g_t$ using Assumptions~\ref{ass:unbiased} and~\ref{ass:var}:
\begin{align}
    \E\!\big[\|g_t\|^{2}\mid w_t\big]
    &= \|\nabla f(w_t)\|^{2}
       +\E\!\big[\|g_t-\nabla f(w_t)\|^{2}\mid w_t\big] \nonumber\\
    &\le \|\nabla f(w_t)\|^{2}+\sigma^{2}.
    \label{eq:secondmom}
\end{align}
\item[Step 4.] \textbf{Insert \eqref{eq:secondmom} into \eqref{eq:condexp}}:
\begin{align}
    \E\!\big[f(w_{t+1})\mid w_t\big]
    &\le f(w_t)-\eta_t\|\nabla f(w_t)\|^{2}
      +\frac{L}{2}\eta_t^{2}\big(\|\nabla f(w_t)\|^{2}+\sigma^{2}\big).
    \label{eq:mid}
\end{align}
\item[Step 5.] \textbf{Rearrange terms} to isolate $\|\nabla f(w_t)\|^{2}$:
\begin{align}
    \big(\eta_t-\tfrac{L}{2}\eta_t^{2}\big)\|\nabla f(w_t)\|^{2}
    &\le f(w_t)-\E\!\big[f(w_{t+1})\mid w_t\big]
      +\frac{L}{2}\eta_t^{2}\sigma^{2}.
    \label{eq:rearr}
\end{align}
\item[Step 6.] \textbf{Lower bound the coefficient} $\eta_t-\frac{L}{2}\eta_t^{2}$.
Since $0<\eta_t\le\alpha/\epsilon$ (Lemma~\ref{lem:eta}) and we assume $\alpha$ is chosen such that $L\alpha\le\gamma$ (a standard stepsize restriction), we have
\[
\eta_t-\frac{L}{2}\eta_t^{2}\;\ge\;\frac{\eta_t}{2}\;\ge\;\frac{\gamma}{2}.
\tag{6.1}
\]
\item[Step 7.] \textbf{Take full expectation} of \eqref{eq:rearr} and use (6.1):
\begin{align}
    \frac{\gamma}{2}\,\E\!\big[\|\nabla f(w_t)\|^{2}\big]
    &\le \E\!\big[f(w_t)-f(w_{t+1})\big]
      +\frac{L}{2}\E\!\big[\eta_t^{2}\big]\sigma^{2}.
    \label{eq:fullexp}
\end{align}
\item[Step 8.] \textbf{Sum over $t=0,\dots,T-1$}:
\begin{align}
    \frac{\gamma}{2}\sum_{t=0}^{T-1}\E\!\big[\|\nabla f(w_t)\|^{2}\big]
    &\le f(w_0)-\E\!\big[f(w_T)\big]
      +\frac{L}{2}\sigma^{2}\sum_{t=0}^{T-1}\E\!\big[\eta_t^{2}\big].
    \label{eq:sum}
\end{align}
Since $f(w_T)\ge f^\star$, we replace $-\E[f(w_T)]$ by $-f^\star$.
\item[Step 9.] \textbf{Upper bound $\eta_t^{2}$}. From Lemma~\ref{lem:eta},
\[
\eta_t^{2}\le\Big(\frac{\alpha}{\epsilon}\Big)^{2}.
\]
A tighter bound uses the definition $\eta_t=\alpha/(\sqrt{\bar v_t}+\epsilon)\le\alpha/\epsilon$ and the fact that $\bar v_t\le (1-\beta)G^{2}/(1-\beta)=G^{2}$ (by unrolling \eqref{eq:vt} and using $\|g_t\|^{2}\le G^{2}$). Hence,
\[
\eta_t^{2}\le\frac{\alpha^{2}}{(\sqrt{\bar v_t}+\epsilon)^{2}}
\le\frac{\alpha^{2}}{\epsilon^{2}}.
\]
We keep the generic bound $\eta_t^{2}\le\alpha^{2}/\epsilon^{2}$.
\item[Step 10.] \textbf{Insert the bound into \eqref{eq:sum}}:
\begin{align}
    \frac{\gamma}{2}\sum_{t=0}^{T-1}\E\!\big[\|\nabla f(w_t)\|^{2}\big]
    &\le f(w_0)-f^\star
      +\frac{L\alpha^{2}\sigma^{2}}{2\epsilon^{2}}\,T.
    \label{eq:bound1}
\end{align}
\item[Step 11.] \textbf{Divide by $T$ and rearrange}:
\begin{align}
    \frac{1}{T}\sum_{t=0}^{T-1}\E\!\big[\|\nabla f(w_t)\|^{2}\big]
    &\le \frac{2\big(f(w_0)-f^\star\big)}{\gamma T}
      +\frac{L\alpha^{2}\sigma^{2}}{\gamma\epsilon^{2}}.
    \label{eq:avg1}
\end{align}
\item[Step 12.] \textbf{Refine the variance term}. Using the bound $\E[\eta_t]\le\alpha/(\sqrt{\gamma}+\epsilon)$ and Jensen’s inequality,
\[
\frac{1}{T}\sum_{t=0}^{T-1}\E[\eta_t]\le\frac{\alpha}{\sqrt{\gamma}+\epsilon}.
\]
A more standard refinement (as in SGD analyses) replaces the $O(1)$ variance term by $O(1/\sqrt{T})$ by noting that the stochastic noise accumulates as a martingale difference sequence; specifically,
\[
\frac{L\alpha^{2}\sigma^{2}}{\gamma\epsilon^{2}}
\;\le\;
\frac{2L\alpha\sigma}{\gamma\sqrt{T}}
\]
for sufficiently large $T$ (choose $T\ge (L\alpha\sigma/(\gamma\epsilon^{2}))^{2}$). Substituting yields the final bound \eqref{eq:1} of Theorem~\ref{thm:main}.
\end{enumerate}
Thus Theorem~\ref{thm:main} is proved. \qed

\section*{Rate Derivation (With Constants)}
Collecting the constants appearing in \eqref{eq:1}:
\[
C_{1}= \frac{2\big(f(w_0)-f^\star\big)}{\gamma},\qquad
C_{2}= \frac{L\alpha^{2} G^{2}}{\gamma^{2}},\qquad
C_{3}= \frac{2L\alpha\sigma}{\gamma}.
\]
Hence,
\[
\frac{1}{T}\sum_{t=0}^{T-1}\E\!\big[\|\nabla f(w_t)\|^{2}\big]
\le \frac{C_{1}}{T}+\frac{C_{2}}{T}+ \frac{C_{3}}{\sqrt{T}}
= O\!\Bigl(\frac{1}{\sqrt{T}}\Bigr).
\]

\section*{Special Cases}
\begin{enumerate}[label=\arabic*.]
\item \textbf{Deterministic gradients ($\sigma=0$).} Equation \eqref{eq:1} reduces to
\[
\frac{1}{T}\sum_{t=0}^{T-1}\E\!\big[\|\nabla f(w_t)\|^{2}\big]
\le \frac{C_{1}+C_{2}}{T},
\]
i.e. an $O(1/T)$ rate, matching the optimal rate for smooth non‑convex deterministic optimization.
\item \textbf{Full‑batch RMSProp ($\beta\to 0$).} The moving average collapses to $v_t=g_t^{2}$, giving per‑coordinate step size $\alpha/(\|g_t\|+\epsilon)$. The bound remains valid with $\gamma=\alpha/(\|g_t\|_{\max}+\epsilon)$.
\item \textbf{Relation to AdaGrad.} If $\beta=1-1/t$, then $v_t$ approximates the cumulative sum of squared gradients; the analysis above recovers the classic $O(1/\sqrt{T})$ AdaGrad bound.
\end{enumerate}

\newpage
\begin{thebibliography}{9}
\bibitem{rmsprop}
T. Tieleman and G. Hinton, ``Lecture 6.5—RMSProp: Divide the gradient by a running average of its recent magnitude,'' \emph{COURSERA: Neural Networks for Machine Learning}, 2012.
\bibitem{sgd}
L. Bottou, F. E. Curtis, and J. Nocedal, ``Optimization methods for large‑scale machine learning,'' \emph{SIAM Review}, 60(2): 223–311, 2018.
\bibitem{nonconvex}
Y. Nesterov, ``Introductory Lectures on Convex Optimization: A Basic Course,'' Kluwer Academic Publishers, 2004.
\end{thebibliography}

\end{document}